\textbf{Procedure.} Same computation for Experiment~B (drone on the ground, rotors spinning, held by hand),
window $t\in[10,33]$~s ($N=1778$), discarding the spin-up and the grabbing of the drone ($t<9$~s).

\begin{table}[H]\centering\scriptsize\setlength{\tabcolsep}{4pt}
\caption{Experiment B: mean and covariance of $z$ (units as in Table~\ref{tab:statsA}).}\label{tab:statsB}
\begin{tabular}{l r rrr rrr}\toprule
 & mean & \multicolumn{6}{c}{covariance $\Sigma$}\\\cmidrule(lr){3-8}
 & & $a_x$ & $a_y$ & $a_z$ & $p$ & $q$ & $r$\\\midrule
$a_x$ & $\phantom{-}0.471$ & $2.86$e-2 & $2.4$e-2 & $-2.4$e-3 & $-5.3$e-4 & $-2.3$e-3 & $-3.8$e-3\\
$a_y$ & $\phantom{-}0.601$ & & $0.606$ & $2.7$e-2 & $4.9$e-3 & $7.2$e-3 & $7.5$e-3\\
$a_z$ & $-9.874$ & & & $2.93$e-2 & $-3.3$e-4 & $1.3$e-3 & $8.5$e-4\\
$p$ & $-3.3$e-3 & & & & $8.25$e-3 & $2.0$e-3 & $6.7$e-3\\
$q$ & $\phantom{-}2.9$e-3 & & & & & $3.18$e-3 & $3.5$e-3\\
$r$ & $-8.7$e-3 & & & & & & $1.08$e-2\\\bottomrule
\end{tabular}
\end{table}

\textbf{Standard deviations:} $\sigma_a=[0.17,\,0.78,\,0.17]$~m/s$^2$, $\sigma_{p,q,r}=[0.091,\,0.056,\,0.104]$~rad/s
($\approx5.2,3.2,6.0^\circ$/s). Relative to A the accelerometer noise grows by $5\times$ ($a_x,a_z$) and $80\times$ ($a_y$), and
the gyro noise by $\approx125\times$ for $p$ ($7.3\times10^{-4}\to0.091$), $77\times$ for $q$ and $146\times$ for $r$. Notable correlations:
$\rho_{pr}=0.71$, $\rho_{qr}=0.60$, $\rho_{pq}=0.39$ (gyros) and $\rho_{a_xa_y}=0.19$, $\rho_{a_xq}=-0.24$.

\textbf{New sources of disturbance.} (i) \emph{Rotor/frame vibration}: four propellers turning at several thousand rpm excite the
airframe and the (non-isolated) IMU at the rotor frequency and its harmonics; since the navdata rate is $\approx77$~Hz, anything above $\approx38$~Hz is
aliased into the measurement and appears as broadband noise on both the accelerometers and the gyros. (ii) \emph{Unbalance and differences of thrust/torque} between rotors
generate real body rates and forces; they appear coupled across axes, which explains the strong $p$--$r$ and $q$--$r$ correlations
(the same torque disturbance excites several axes) and the non-zero off-diagonal acceleration terms. (iii) \emph{The hand holding the drone}: the
tremor and the small reaction movements of the hand give low-frequency real accelerations and rotations; the large $\sigma_{a_y}=0.78$ and the mean $\bar a_y=0.60$ (versus $-0.37$ in A) show
that the airframe is tilted/pushed sideways ($\approx3.5^\circ$ roll, with the roll swinging by $\sigma\approx4.5^\circ$ in the raw inclinometer, Section~3.4) and no longer at rest. (iv) \emph{Aerodynamic disturbances}: the downwash is reflected by the ground/the
hand (ground effect), producing a turbulent and non-stationary flow on the frame. As a result the gyro biases of A ($\sim10^{-3}$~rad/s)
become negligible compared with the noise ($\sim10^{-1}$~rad/s), and the accelerometer mean is no longer the pure gravity vector ($\bar a_x=0.47$, $\|\bar a\|=9.91$~m/s$^2$).
Note that the noise is no longer white and stationary (the std over 3~s blocks varies between $0.09$ and $0.14$~rad/s for $r$), so a single
covariance is only an average description.
